\section{Problem Analysis: Why Agentic Lifecycles Need a Verification Spine}
\label{sec:motivation}

The emerging generation of AI-native software lifecycles---typified by AWS AI-DLC~\cite{awsaidlc}, GitHub Spec Kit~\cite{githubspeckit}, and the Kiro spec-driven IDE~\cite{awskiro}---share a common and genuinely attractive premise: that natural language (NL) specifications, rather than handwritten code, should be the primary artefact a developer authors, with AI agents responsible for translating those specifications into running implementations. This premise is supported by a strong empirical trend. Repository-level coding agents now resolve real issues against full codebases~\cite{jimenez2024swebench,yang2024sweagent,zhang2024autocoderover}, and the broader position literature anticipates a corresponding restructuring of the software development lifecycle around agentic generation~\cite{ozkaya2023llmse,liu2024llmagentse}. The thesis of this article is that the premise is correct but its current realizations are structurally unsound: they elevate the specification to the role of source of truth while leaving the specification, and the loop that consumes it, in a form that cannot be machine-checked. This section analyses precisely why, isolates three named defects shared across these lifecycles, frames the economic inversion that makes the defects consequential, and motivates the reference architecture developed in the remainder of the article.

\subsection{Three structural defects of current spec-driven lifecycles}
\label{subsec:three-defects}

Examined as a control system rather than as a developer experience, the AI-DLC, Spec Kit, and Kiro lifecycles exhibit three defects that compound one another. They are stated here as named problems because the architecture in \cref{sec:architecture} is organised as a direct response to each.

\paragraph{Defect 1: the specification is natural language, so spec-to-code is an unverifiable compilation.}
In each of these lifecycles the durable input is a NL document---a prompt, a ``spec'' file, or a requirements note. The agent's job is to ``compile'' that document into code. But because the source language of this compilation is NL, there is no formal relation between source and target that any tool can check. A conventional compiler can be tested, and in principle verified, because both its input and output have a defined semantics; a NL-to-code agent has, on the input side, only the diffuse and underdetermined meaning of a human sentence. The consequence is that the central transformation of the lifecycle is, by construction, outside the reach of mechanical verification. This is qualitatively different from the situation in classical program synthesis, where the input is itself a formal contract and the synthesized program can be certified against it~\cite{manna1980deductive,solarlezama2008sketching}.

\paragraph{Defect 2: the loop is open---generate then eyeball.}
Having produced code, these lifecycles return it to a human for inspection, or run a finite, human-selected test suite. There is no \emph{machine gate} that admits an implementation only when it provably meets the intent. We call this the \emph{open loop}: the feedback path that should close ``code back onto specification'' is completed, if at all, by human judgment or by sampled tests. The reliability literature has documented in detail how unreliable that closure is. Functional-correctness estimates from test suites are systematically optimistic because the suites are too weak to catch plausible-but-wrong solutions, and strengthening them lowers reported pass rates substantially~\cite{liu2023evalplus,chen2021codex}. LLM-generated code is moreover prone to confident fabrication---non-existent APIs, intent conflicts, subtle off-by-one and boundary errors---that survives review precisely because the surrounding code is well-formed~\cite{zhang2025hallucination}. An open loop over a fallible generator is not a quality process; it is a presentation layer.

\paragraph{Defect 3: where verification exists, it is circular---self-grading.}
The few points at which these lifecycles do attempt automated acceptance rely on tests, and those tests are typically generated by the same model, in the same session, against the same NL understanding, that produced the code. The generator thereby writes its own oracle. Because the test author and the implementation author share a model and a training distribution, an error in the model's understanding of the intent is replicated identically in code and in oracle, and the two agree---wrongly---without contradiction. The hazard has been observed directly in LLM-driven oracle and assertion generation, where the test and the code drift together away from the true intent~\cite{liu2026oracle,watson2020atlas}. We call this the \emph{self-grading} or circularity problem. It is the deepest of the three defects because it can make verification appear to succeed while certifying nothing.

These defects are not independent failures but a single missing component viewed from three angles. The missing component is a \emph{verification spine}: a machine-checkable specification that is (a) related to code by a checkable relation, closing Defect~1; (b) the gate of a closed loop, closing Defect~2; and (c) owned independently of the implementation agent, closing Defect~3.

\subsection{A motivating example: weakening ``sort'' to ``is a permutation''}
\label{subsec:motivating-example}

The defects above are easiest to see in a concrete failure. Consider a developer whose NL intent is ``sort the array ascending.'' Suppose the agent is asked, as in current spec-driven flows, to record this intent as a formal contract and then to satisfy it. A faithful postcondition for an in-place sort over \texttt{int[] a} has two conjuncts: the result is sorted, and the result is a permutation of the input. \Cref{lst:weak-sort} shows the faithful contract beside a contract the agent might propose if it is permitted to author the oracle it will later be graded against.

\begin{lstlisting}[caption={A faithful sort contract (top) and a silently weakened contract (bottom). The weakened contract drops the orderedness conjunct; the identity function satisfies it.},label={lst:weak-sort}]
//@ ensures (\forall int i; 0 <= i && i < a.length-1; a[i] <= a[i+1]);
//@ ensures isPerm(\old(a), \result); // result is a permutation of the input
int[] sort(int[] a) { ... }

//@ ensures isPerm(\old(a), \result); // "is a permutation" -- orderedness dropped
//@ assignable \nothing;
int[] sort(int[] a) { return a; }     // identity satisfies the weakened contract
\end{lstlisting}

Here \texttt{isPerm} denotes a user-defined \texttt{pure} model predicate, since JML provides no built-in permutation operator. The weakened contract retains only the permutation conjunct. It is still verifiable; indeed it is verifiable by the identity function, which returns the input unchanged (a permutation of itself). Under a circular pipeline the same model that drafted this contract then writes \texttt{return a}, the verifier reports success, and the lifecycle declares the task done. Nothing in the open loop catches the omission: the code provably satisfies the contract, the contract is plausible NL-adjacent text, and a human glancing at a green check has no reason to object. This is the self-grading defect (Defect~3) producing a spec-weakening pathology that no amount of code-side verification can detect, because the code is correct \emph{relative to the weakened contract}. The error lives entirely in the specification, and is invisible unless the specification is itself adversarially scrutinized against the original intent---a mechanism we develop as adversarial elaboration in \cref{sec:architecture}.

A symmetric pathology arises in brownfield code. Consider a pre-existing method whose Javadoc states ``returns the index of \texttt{key}, or $-1$ if absent,'' but whose body, through an off-by-one in a refactor, returns $0$ rather than $-1$ on the empty array. Here the code and the documentation \emph{disagree}, and that disagreement is exactly the bug. Inferring a contract from the code alone~\cite{ernst2007daikon,edmonds_article3} would faithfully encode \texttt{result == 0} for the empty case and verify it green---characterising the bug rather than finding it. The bug surfaces only when the code-derived clause is reconciled against the documentation-derived intent and the conflict is flagged. Both examples make the same point: verifying code against the wrong specification is not merely unhelpful, it is actively misleading, and a wrong specification verified against is worse than no specification at all.

\subsection{The economic inversion: code is disposable, the specification is durable}
\label{subsec:economic-inversion}

The defects of \cref{subsec:three-defects} have always existed in informal practice; what makes them urgent now is an inversion in the cost structure of software. When an implementation can be regenerated by an agent in seconds and at negligible marginal cost, the implementation ceases to be the scarce, durable artefact it was for the preceding half-century of software engineering. Code becomes a \emph{disposable build output}---a derivative that can be discarded and re-synthesized whenever requirements shift, the surrounding library set changes, or a faster model becomes available. The artefact that retains value across regenerations is the precise statement of \emph{what the code must do}: the specification. This is the economic inversion, and it has a direct architectural corollary. A lifecycle that treats fluid, regenerable code as its source of truth and an unchecked NL note as a transient input has the durability gradient backwards. The lifecycle must instead be reorganised so that the machine-checkable specification is the persistent source of truth and code is the regenerable projection of it.

This inversion is what converts the three defects from chronic annoyances into load-bearing failures. If code were the durable artefact, an occasional spec-weakening or a stale comment would be a local blemish to be patched in the implementation. Once the specification is the durable artefact---the thing every regeneration is built against and graded by---a defect in the specification is inherited by every future implementation. The cost of an unverifiable, self-graded, open-loop specification is therefore no longer paid once; it compounds. The classical disciplines that anticipated this regime---Hoare logic and weakest-precondition reasoning~\cite{hoare1969axiomatic,dijkstra1976discipline}, Design by Contract~\cite{meyer1992dbc,meyer1997oosc}, and refinement calculus~\cite{back1998refinement,morgan1994programming}---already located correctness in the relation between a specification and an implementation rather than in the implementation alone. The contribution of this article is to make that relation the operational pivot of an AI-native lifecycle, using a contract language and verifier that are already mature for Java: JML~\cite{leavens2006jml} discharged by OpenJML extended static checking~\cite{cok2011openjml,cok2014openjml} over the Z3 SMT solver~\cite{demoura2008z3}.

\subsection{The ratification gap: the central tension of the intent layer}
\label{subsec:ratification-gap}

Reorganising the lifecycle around a machine-checkable contract immediately raises a tension that, left unresolved, would reintroduce every defect it was meant to remove. We call it the \emph{ratification gap}. The argument is short and unforgiving. If the formal contract is the oracle against which all code is graded, then the contract must be correct, and---because correctness is relative to human intent, which no tool can divine---the contract must ultimately be owned and confirmed by a human. Yet the human chose NL in the first place precisely because they cannot, or will not, author and read formal contracts. Asking that same human to confirm a page of quantified JML they cannot decode does not establish human authority over the oracle; it manufactures the \emph{appearance} of authority while the human rubber-stamps text they do not understand. A pipeline built on fake ratification is no better grounded than the open loop it replaces: the self-grading agent is simply relabeled as ``human-approved.''

The ratification gap is the central design problem of the intent layer, and the architecture in \cref{sec:architecture} is structured around refusing the false resolutions of it. It does not ask the human to read formal logic; instead it ratifies \emph{behaviour, not logic}, rendering each contract as concrete, checkable examples---``given \texttt{x = -1}, this contract says the result is \texttt{0}; is that correct?''---so that the human exercises genuine authority over observable behaviour while the formal text remains the machine's concern. It surfaces, Socratically, only the boundary decisions the agent had to resolve (null and empty inputs, overflow, ordering, duplicates, concurrency, error-versus-exception), rather than burying the human in the full contract. And it enforces an \emph{independent-authority constraint} so that the contract the code is graded against is never authored by the agent that writes the code. \Cref{tab:defect-mechanism} summarises the correspondence between the defects analysed here and the mechanisms that answer them, each defined in the architecture that follows.

\begin{table}[t]
\centering
\caption{The three defects of current agentic lifecycles, the named problems they give rise to, and the architectural mechanism (\cref{sec:architecture}) that responds to each.}
\label{tab:defect-mechanism}
\begin{tabular}{@{}lll@{}}
\toprule
Defect & Named problem & Responding mechanism \\
\midrule
NL spec, unverifiable compilation & Unverifiable intent & Formal contract as pivot; verify all inputs \\
Generate-then-eyeball & Open loop & CEGIS-shaped synthesis--verification loop \\
Same model writes code + oracle & Self-grading / circularity & Independent-authority constraint \\
Human cannot read formal text & Ratification gap & Ratify behaviour, not logic (pinned oracles) \\
Spec weaker than intent & Spec weakening & Adversarial elaboration; vacuity checks \\
Code-derived spec only describes & Near-circular characterisation & Multi-source intent triangulation \\
\bottomrule
\end{tabular}
\end{table}

\subsection{Why inference makes the reorganisation feasible}
\label{subsec:inference-bootstrap}

A final objection must be met before the architecture can be credible: if formal contracts are so expensive to author that developers reach for NL, does reorganising the lifecycle around them not simply relocate the cost rather than remove it? The answer, and the reason this programme is positioned to attempt the reorganisation, is that specification \emph{inference} removes the authoring bootstrap. The author's prior work in this programme has established the necessary instrument and its lifecycle fit: inferred JML specifications materially improve LLM-based unit-test generation~\cite{edmonds_article1}; those specifications can be embedded into compiled artefacts and OpenAPI documents at full roundtrip fidelity~\cite{edmonds_article2}; heuristic inference admits compositional refinement, with a second weakest-precondition pass adding roughly $18{,}854$ additional well-formed preconditions across the study corpus~\cite{edmonds_article3}; and the inference pipeline integrates into CI/CD with diff-only analysis, content-hash caching, and fail-open semantics~\cite{edmonds_article4}. Inference supplies a draft contract automatically, so the human's task collapses from \emph{authoring} formal text to \emph{ratifying} behaviour---exactly the operation the ratification gap demands be cheap. The compositional result is especially consequential for the lifecycle, because it means a change to one method's contract recomputes the proof obligations of its callers without re-running them, turning internal refactors and dependency version bumps into the same first-class, statically answerable question: which of my call sites does this change break? \Cref{sec:architecture} develops the layered model that turns these capabilities into a verification spine, and \cref{sec:instantiation} reuses the inferrer, the OpenJML fork, and the Z3 oracle as trusted components of a proof-of-concept instantiation, so that the verifier never grades itself.
